\textbf{1. Kiến thức cần nhớ}
\begin{itemize}
\item \textbf{Phép thử ngẫu nhiên} là một hành động, thực nghiệm mà kết quả của nó không thể biết trước, nhưng ta biết trước được tập hợp tất cả các kết quả có thể xảy ra.
\item \textbf{Không gian mẫu} là tập hợp tất cả các kết quả có thể xảy ra của phép thử, kí hiệu là $\Omega$. Số phần tử của không gian mẫu kí hiệu là $n(\Omega)$.
\item Các công cụ liệt kê không gian mẫu: Liệt kê trực tiếp, dùng bảng $2$ chiều, dùng sơ đồ hình cây.
\end{itemize}

\textbf{2. Phương pháp giải các dạng}

%% DANG 9D8?25-1
\textbf{Dạng 1. Mô tả phép thử, xác định không gian mẫu}
\begin{enumerate}
\item Xác định phép thử được thực hiện qua mấy công đoạn (ví dụ: gieo $2$ đồng xu, rút $2$ thẻ).
\item Liệt kê đầy đủ các phần tử của tập hợp $\Omega$ dưới dạng các cặp hoặc bộ số $(x; y)$.
\item Đếm số phần tử $n(\Omega)$.
\end{enumerate}

%% DANG 9D8?25-2
\textbf{Dạng 2. Liệt kê kết quả bằng bảng, sơ đồ cây}
\begin{enumerate}
\item \textbf{Dùng bảng hai chiều}: Thích hợp khi phép thử gồm hai công đoạn độc lập (gieo $2$ con xúc xắc: bảng $6 \times 6 = 36$ ô).
\item \textbf{Dùng sơ đồ hình cây}: Vẽ các nhánh từ gốc: mỗi nhánh ở tầng $1$ là một kết quả của lần $1$, từ mỗi nhánh đó rẽ tiếp các nhánh ở tầng $2$ cho lần $2$. Số nhánh tận cùng là số kết quả có thể xảy ra.
\end{enumerate}
